% Standalone problem document, generated from the adjacent README.md.
% Compile with XeLaTeX. Mathematical content is maintained in Markdown.
\documentclass[11pt,a4paper]{article}
\usepackage[fontsize=10.5pt]{scrextend}
\usepackage[margin=22mm,headheight=15pt,headsep=9mm,footskip=11mm]{geometry}
\usepackage{fontspec}
\setmainfont{texgyrepagella}[Extension=.otf,UprightFont=*-regular,BoldFont=*-bold,ItalicFont=*-italic,BoldItalicFont=*-bolditalic]
\setsansfont{texgyreheros}[Extension=.otf,UprightFont=*-regular,BoldFont=*-bold,ItalicFont=*-italic,BoldItalicFont=*-bolditalic]
\setmonofont{texgyrecursor}[Extension=.otf,UprightFont=*-regular,BoldFont=*-bold,ItalicFont=*-italic,BoldItalicFont=*-bolditalic,Scale=MatchLowercase]
\usepackage{amsmath,amssymb}
\usepackage{unicode-math}
\setmathfont{texgyrepagella-math.otf}
\usepackage{xcolor,microtype,enumitem,needspace,fancyhdr}
\usepackage{bookmark}
\definecolor{ink}{HTML}{17344B}
\definecolor{accent}{HTML}{197D80}
\definecolor{muted}{HTML}{596773}
\hypersetup{colorlinks=true,linkcolor=accent,urlcolor=accent,pdftitle={KE-04 - Strict
interlacing across block Lanczos
iterations},pdfauthor={Open Problems in Numerical Linear Algebra}}
\urlstyle{same}
\setlength{\parindent}{0pt}
\setlength{\parskip}{5pt plus 1pt minus 1pt}
\linespread{1.04}
\setlength{\emergencystretch}{3em}
\widowpenalty=10000
\clubpenalty=10000
\displaywidowpenalty=10000
\allowdisplaybreaks[1]
\setlist{nosep,leftmargin=1.5em}
\providecommand{\tightlist}{\setlength{\itemsep}{0pt}\setlength{\parskip}{0pt}}
\providecommand{\pandocbounded}[1]{#1}
\pagestyle{fancy}
\fancyhf{}
\fancyhead[L]{\sffamily\footnotesize\color{muted}OPEN PROBLEMS IN NUMERICAL LINEAR ALGEBRA}
\fancyhead[R]{\sffamily\bfseries\footnotesize\color{accent}KE-04}
\fancyfoot[L]{\parbox[b]{0.9\textwidth}{\sffamily\fontsize{8}{10}\selectfont\color{muted}Editorial ratings; a bounded search cannot certify that a problem remains unsolved.\\Literature check: 2026-09-13}}
\fancyfoot[R]{\sffamily\footnotesize\color{muted}\thepage}
\renewcommand{\headrulewidth}{0.3pt}
\renewcommand{\headrule}{\hbox to\headwidth{\color{accent}\leaders\hrule height \headrulewidth\hfill}}
\makeatletter
\renewcommand{\section}{\Needspace{5\baselineskip}\@startsection{section}{1}{0pt}{12pt plus 2pt minus 2pt}{4pt}{\sffamily\bfseries\large\color{ink}}}
\renewcommand{\subsection}{\Needspace{4\baselineskip}\@startsection{subsection}{2}{0pt}{9pt plus 1pt minus 1pt}{3pt}{\sffamily\bfseries\normalsize\color{ink}}}
\makeatother
\setcounter{secnumdepth}{0}
\begin{document}
{\sffamily\small\color{muted}Eigenvalues and inverse problems\par}
\vspace{3pt}
{\raggedright\hyphenpenalty=10000\exhyphenpenalty=10000\sffamily\bfseries\fontsize{21}{24}\selectfont\color{ink}Strict
interlacing across block Lanczos iterations\par}
\vspace{7pt}
{\sffamily\small\textbf{Difficulty:} challenging\quad\textbf{Importance:} interesting
to the community\par}
{\sffamily\small\bfseries\color{accent}Status: Lean verified\par}
\vspace{4pt}
{\color{accent}\hrule height 0.8pt}
\vspace{6pt}
\textbf{Rating rationale:} Challenging reflects a stronger strict
interlacing law than general compression interlacing supplies; community
impact is spectral information across block Krylov iterations.

\subsection{Lean proof and verification evidence ---
2026-09-13}\label{lean-proof-and-verification-evidence-2026-09-13}

\textbf{Formalization: George Stepaniants}, Department of Computing and
Mathematical Sciences, California Institute of Technology, Pasadena,
California, USA. Matthew J. Colbrook retains credit for the original
mathematical proof; D. Šimonová and P. Tichý retain credit for the
conjecture. The formalization used substantial AI assistance.

The
\href{https://github.com/sgstepaniants/OpenProblemsInNLA/tree/40b0bf52e73e776e7769f0f12dbbda7cd9fff183/eigenvalues-and-inverse-problems/KE-04/lean}{immutable
Lean source} proves \textit{NLA.KE04.blockLanczosConjecture}, the
complete original largest-full-iteration statement. The proof derives it
from the stronger full-prefix theorem. All 24
\href{https://github.com/ajt60gaibb/OpenProblemsInNLA/blob/main/eigenvalues-and-inverse-problems/KE-04/lean/comparator.json}{registered
theorem declarations} are checked. The proof retains every dimension,
real symmetric matrix, full-rank starting block, allowed iteration pair,
strict interval index and independently chosen orthonormal Krylov basis.
Eigenvalue multiplicities are retained, and no separated-endpoint or
compatible-basis assumption is added. See the
\href{https://github.com/ajt60gaibb/OpenProblemsInNLA/blob/main/eigenvalues-and-inverse-problems/KE-04/lean/SourceCorrespondence.md}{historical
statement-stage correspondence} and
\href{https://github.com/ajt60gaibb/OpenProblemsInNLA/blob/main/eigenvalues-and-inverse-problems/KE-04/lean/PROOF-MAP.md}{proof
map}. The correspondence is preserved from before proof implementation;
its pending-gate descriptions are historical and superseded by the
completed verification below.

Lean is pinned to \textbf{4.33.1}, Mathlib to
\textit{0df444a360eaa60ab8c11dca51a86af692955474}, and LeanCert to
\textit{621a43d7cf21f87872392a01e874f2f1dbddc926}. The argument uses
exact finite-dimensional algebra; LeanCert checks the actual proof
declarations in kernel mode. The
\href{https://github.com/ajt60gaibb/OpenProblemsInNLA/blob/main/eigenvalues-and-inverse-problems/KE-04/lean/lake-manifest.json}{full
dependency manifest} fixes all ten packages.

This campaign ran the real Ubuntu verifier:
\href{https://github.com/sgstepaniants/OpenProblemsInNLA/actions/runs/34759746409/job/103730400358}{workflow
34759746409, KE-04 job 103730400358}. Comparator accepted the 24
independently reviewed statements and proofs, Lean's default kernel
accepted the exported proofs, and the transitive axiom reports permit
only \textit{propext}, \textit{Classical.choice}, and
\textit{Quot.sound}. The actual sandbox and rejection controls passed.
The
\href{https://github.com/ajt60gaibb/OpenProblemsInNLA/blob/main/eigenvalues-and-inverse-problems/KE-04/lean/verification/linux-2026-09-13/README.md}{retained
logs and independent operational audit} bind the run to the immutable
source revision. The proof also passed two independent final
mathematical reviews:
\href{https://github.com/ajt60gaibb/OpenProblemsInNLA/blob/main/eigenvalues-and-inverse-problems/KE-04/lean/reviews/final-referee-1.md}{referee
1} and
\href{https://github.com/ajt60gaibb/OpenProblemsInNLA/blob/main/eigenvalues-and-inverse-problems/KE-04/lean/reviews/final-referee-2.md}{referee
2}. These are AI-agent reviews, not external human peer review.

From a clean checkout, reproduce on non-root Linux with the
\href{https://github.com/ajt60gaibb/OpenProblemsInNLA/blob/main/tools/lean/HARNESS.md}{documented
prerequisites}:

\begin{verbatim}
tools/lean/bootstrap.sh /tmp/nla-ke04-check
tools/lean/verify.sh \
  eigenvalues-and-inverse-problems/KE-04/lean /tmp/nla-ke04-check
\end{verbatim}

\subsection{Resolution --- 2026-09-11}\label{resolution-2026-09-11}

\textbf{Affirmative resolution by Matthew J. Colbrook} (Department of
Applied Mathematics and Theoretical Physics, University of Cambridge).
See the
\href{https://github.com/ajt60gaibb/OpenProblemsInNLA/blob/main/eigenvalues-and-inverse-problems/KE-04/solution.md}{complete
manuscript}, \textbf{Theorem KE-04, sections 1--3}
(\href{https://github.com/ajt60gaibb/OpenProblemsInNLA/blob/main/eigenvalues-and-inverse-problems/KE-04/solution.pdf}{PDF}
·
\href{https://github.com/ajt60gaibb/OpenProblemsInNLA/blob/main/eigenvalues-and-inverse-problems/KE-04/solution.tex}{LaTeX}),
prepared 11 September 2026.

Strict interval occupancy holds for every allowed pair of block Lanczos
iterations and every indicated index, in exact arithmetic before the
first loss of full block dimension. The quadratic-polynomial argument
includes multiplicities and excludes coincident interval endpoints in
the stated range.

The original proof draft was generated in a ChatGPT conversation. A
separate Codex agent independently verified the full proof and its match
to the exact target on 11 September 2026:
\href{https://github.com/ajt60gaibb/OpenProblemsInNLA/blob/main/references/colbrook-2026-09-11/verification/reviews/KE-04-review.md}{detailed
PASS review}. The review records a hash of the unchanged proof text.
This is independent agent verification, not external human peer review
or formal certification.
\href{https://github.com/ajt60gaibb/OpenProblemsInNLA/blob/main/references/colbrook-2026-09-11/README.md}{The
submission history and diagnostic record} preserve the initial solution
claim. The ratings above are historical, and the earlier literature
checks below are retained.

\subsection{Original problem
statement}\label{original-problem-statement}

Let \(A\in\mathbb R^{n\times n}\) be symmetric and
\(V\in\mathbb R^{n\times p}\) have full column rank. Write

\[\mathcal K_j=\mathop{\mathrm{range}}\nolimits[V,AV,\ldots,A^{j-1}V],\]

and let \(s\) be the largest integer with \(\dim\mathcal K_s=sp\). For
\(1\le j\le s\), choose an orthonormal basis \(Q_j\) of \(\mathcal K_j\)
and order the eigenvalues of \(T_j=Q_j^TAQ_j\) as
\(\theta_1^{(j)}\le\cdots\le\theta_{jp}^{(j)}\), with multiplicity.

For every \(1\le k< j\le s\) and \(1\le i\le(k-1)p\), must the open
interval

\[(\theta_i^{(k)},\theta_{i+p}^{(k)})\]

contain an eigenvalue of \(T_j\)? The statement concerns exact
arithmetic before the first loss of full block dimension. Although
Lanczos supplies compatible block tridiagonal representations, the
spectra do not depend on the chosen bases.

This would extend a scalar Lanczos property used to interpret later Ritz
values and distinguish genuine eigenvalue approximation from clusters
created by finite precision.

\newpage
\subsection{References}\label{references}

D. Šimonová and P. Tichý, \emph{On finite precision block Lanczos
computations}, arXiv:2507.16484v1 (22 July 2025), §2.1, unnumbered
conjecture and final discussion
(\href{https://arxiv.org/html/2507.16484v1}{primary manuscript}). J.
Liesen and Z. Strakoš, \emph{Krylov Subspace Methods: Principles and
Analysis}, Oxford 2013, the scalar Lanczos/orthogonal-polynomial
background cited by the source
(\href{https://doi.org/10.1093/acprof:oso/9780199655410.001.0001}{publisher}).

\subsection{Status check --- 2026-09-08}\label{status-check-2026-09-08}

The arXiv abstract/history and full §2.1 were inspected; v1 remains the
only listed version. Searches ``block Lanczos interlacing conjecture'',
``Šimonová Tichý interlacing'', and 2026 proof/counterexample queries
found no resolution. Standard weak Cauchy interlacing is not the
asserted strict interval-occupancy result.

\subsection{Audit update --- 2026-09-10}\label{audit-update-2026-09-10}

Rechecked the explicit conjecture and concluding discussion in the
\href{https://arxiv.org/html/2507.16484v1}{2025 block-Lanczos
manuscript}. Its block-width index ranges agree with this entry.
Searches for later strict block-Lanczos interlacing proofs found no
resolution; ordinary Cauchy interlacing is weaker than the displayed
claim.
\end{document}
